% Solutions of Book 4, chapter 20 (the application layer).
\section*{Chapter \ref{ch:b2:application-layer} --- The Application Layer: DNS, HTTP and TLS}

\begin{solution}{exo:b2:application-layer:1}
\omoutput{code/bachelor-2/20-application-layer/outputs/exo_header.txt}{The header, field by field.}
QR is 1: it is an answer, to the query with identifier \code{2a17}. The server is authoritative (AA), the query asked
for recursion (RD) and the server offers it (RA). RCODE 3 says that the name asked does not exist, which is why there
are no answers; the question is repeated, as in every answer.
\end{solution}

\begin{solution}{exo:b2:application-layer:2}
\omoutput{code/bachelor-2/20-application-layer/outputs/exo_loop.txt}{A pointer that goes back, then one that points
at its own name.}
The name starts at 12 with the label \code{www}, and the pointer at 16 sends the reader back to 12: a reader that
follows every pointer reads \code{www.www.www}\dots\ for ever, and a 6-byte trick stops a careless server. With the
rule of \cref{prop:b2:application-layer:compression}, each pointer lands strictly before the previous start, so a
reader follows fewer than $n/2$ pointers (each one starts a stretch of at least two bytes). Without the rule, a reader
can still stop: on a path without a loop it meets each pointer at most once, and a message of $n$ bytes holds at most
$n/2$ pointers, so a reader that has followed more than $n/2$ has met a loop and may reject the message. The limit of
255 bytes on a name does not suffice alone, since \code{c0 0c} at offset 12 loops without adding a byte.
\end{solution}

\begin{solution}{exo:b2:application-layer:3}
\omoutput{code/bachelor-2/20-application-layer/outputs/exo_aaaa.txt}{The query, twelve bytes a line.}
The header: identifier \code{0b 0e}, flags \code{01 00} (only RD), one question, no records. The name: \code{03}
\code{www}, \code{08} \code{brenmoor}, \code{04} \code{test}, \code{00}; then the type \code{00 1c} (28, AAAA) and the
class \code{00 01} (Internet): 35 bytes.
\end{solution}

\begin{solution}{exo:b2:application-layer:4}
Empty cache: three queries (root, \code{test}, \code{brenmoor.test}). With the referral to \code{brenmoor.test}
cached: one, to the town's server. With the answer cached: none. The machine waits one round trip to its \omterm{def:b2:application-layer:resolver}{resolver},
plus the \omterm{def:b2:application-layer:resolver}{resolver}'s queries: in the three cases 1 + 3, 1 + 1 and 1 round trip, of different lengths since the
\omterm{def:b2:application-layer:resolver}{resolver} is usually much closer than the name servers.
\end{solution}

\begin{solution}{exo:b2:application-layer:5}
\omcode{code/bachelor-2/20-application-layer/exercises/chunked.py}{4}{20}{Decoding a chunked body}
Each chunk is a size line, the data and a line end; the size is read in hexadecimal, and anything after a \code{;} on
the size line (an extension) is ignored. The last chunk has size 0 and is followed by an empty line. Returning the
number of bytes used lets the caller read the next response on the same connection. The test feeds it a body with an
extension and with line ends inside the data, and three cut or lying bodies.
\end{solution}

\begin{solution}{exo:b2:application-layer:6}
\omoutput{code/bachelor-2/20-application-layer/outputs/exo_content_length.txt}{The client, with a timeout of 1 s.}
The client trusts \code{Content-Length} and waits for the two missing bytes. On a \omterm{def:b2:application-layer:persistent}{persistent connection} the server
keeps the connection open, so nothing tells the client that no more will come: without a timeout it waits for ever.
Had the server closed the connection, the client would have seen a body shorter than announced and reported an error.
\end{solution}

\begin{solution}{exo:b2:application-layer:7}
\omcode{code/bachelor-2/20-application-layer/listings/http_server.py}{27}{33}{Keeping the connection, and
\code{HEAD}}
The loop of \cref{lst:b2:application-layer:server} reads the next request on the same connection, and leaves only on
a close asked by the client (or an HTTP/1.0 request, whose default was to close) or when the client closes. For
\code{HEAD}, the headers are those of \code{GET}, including the \code{Content-Length} the body would have had, and the
body is not sent. The test sends \code{GET}, \code{HEAD} and \code{GET} on one connection and checks one connection,
the lengths and the empty body.
\end{solution}

\begin{solution}{exo:b2:application-layer:8}
Cut time into cycles, each starting with a query sent upstream (a miss). The answer is then cached for $T$ seconds, and
every query in that interval is a hit; the first query after it is the next miss. In a Poisson process the number of
queries in an interval of length $T$ has mean $rT$, whatever happened before. A cycle therefore holds one miss and on
average $rT$ hits, and over many cycles the fraction of hits is
\[
\frac{rT}{1 + rT}.
\]
\cref{out:b2:application-layer:cache} agrees to within 0.0005: with $rT = 1$ (one query a minute, a lifetime of a
minute), half the queries are hits.
\end{solution}

\begin{solution}{exo:b2:application-layer:9}
\omcode{code/bachelor-2/20-application-layer/exercises/caching_resolver.py}{8}{39}{A caching iterative \omterm{def:b2:application-layer:resolver}{resolver}}
The cache maps a name and a type to records and an expiry time; referrals are kept under the zone's name, as the
address of its server (the glue). A resolution first looks for the answer, then for the closest enclosing zone whose
server it knows, and starts there instead of at the root. The test runs it on the toy network with a clock it moves by
hand: \code{www} costs 3 queries, then 0; \code{mail} costs 1; after 4000 s \code{www} costs 1 again, and after two
days 3.
\end{solution}

\begin{solution}{exo:b2:application-layer:10}
Packet $i$ (from 0 to 31) belongs to page $i \bmod 4$, leaves at $25 + i$ ms and arrives at $50 + i$. Packet 4,
page 0's second, is lost; packet 7 leaves at 32 ms, so the copy arrives at $32 + 75 = 107$ ms. Each page's last packet
arrives at 78, 79, 80 and 81 ms. With HTTP/2, TCP delivers bytes in order: nothing after packet 4 reaches the
application before 107 ms, so all four pages complete at 107. With HTTP/3, each stream waits only for its own
packets: page 0 completes at 107, pages 1 to 3 at 79, 80 and 81, as without loss. HTTP/1.1 is slowest because each
page waits for the previous one and for its own request to cross the network: a page takes $25 + 7 + 25 = 57$ ms
after its request leaves, and page 0 takes until 107 because of the loss, so the others complete at 164, 221 and
278 ms.
\end{solution}

\begin{solution}{exo:b2:application-layer:11}
\omcode{code/bachelor-2/20-application-layer/listings/tls_echo.py}{5}{26}{A TLS echo pair that checks the server's
name}
The client's context trusts only the given certificate (\code{cafile}), and \code{server_hostname} makes the module
check that the certificate names \code{www.brenmoor.test}; \code{create_default_context} turns both checks on. The
server reads one message and sends it back in capitals: TLS gives it plain bytes, as TCP would.
\end{solution}

\begin{solution}{exo:b2:application-layer:12}
\omoutput[8-16]{code/bachelor-2/20-application-layer/outputs/tls_demo.txt}{Bytes on the wire after the
handshake.}
Each record adds 22 bytes: the 5-byte header, the 1-byte real content type that TLS 1.3 places inside the encrypted
part, and a 16-byte authentication tag that lets the receiver detect any change. A record holds at most 16\,384 bytes,
so 16\,385 bytes take two records and 44 extra bytes. The overhead is $22/23 \approx 96~\%$ of the traffic for a 1-byte
message, $22/122 \approx 18~\%$ for 100 bytes and $22/16\,406 \approx 0.13~\%$ for 16\,384: small messages should be
grouped before they are sent.
\end{solution}

\begin{solution}{pb:b2:application-layer:1}
\textbf{1.} NS: the zone's server; A: the addresses of \code{ns}, \code{www} and \code{mail}; AAAA: the \omterm{def:b2:network-layer:ipv6}{IPv6 address}
of \code{www}; MX: the town's mail server, preference 10; TXT: a line of text; CNAME: \code{library} is another name
of \code{www}.

\textbf{2.} To reach \code{ns.brenmoor.test} a \omterm{def:b2:application-layer:resolver}{resolver} would have to resolve a name inside the zone it is trying to
reach, a circle. The glue record in the parent zone breaks it.

\textbf{3.} The answer repeats the question and adds one record: a pointer to the name (2 bytes), type, class,
lifetime and length (10), the preference (2), the label \code{mail} (5) and a pointer to \code{brenmoor.test} (2):
21 bytes.

\textbf{4.} The CNAME record and the A record of \code{www}, in one answer, so that the \omterm{def:b2:application-layer:resolver}{resolver} needs no second
query.

\textbf{5.} That the name does not exist at all, for any type; RCODE 0 with no answer means that the name exists
without records of that type.

\textbf{6.} A loop would keep the server busy for ever on a message of a few bytes; by
\cref{prop:b2:application-layer:compression}, refusing it makes every read end.

\textbf{7.} 3, then 1 (the referral to \code{brenmoor.test} is cached), then 0.

\textbf{8.} 3: every record has expired, the longest-lived (the referral to \code{test}) after 172\,800 s, two days.

\textbf{9.} $r = 1000/86\,400$ per second and $rT \approx 41.7$, so a fraction $1/(1 + rT) \approx 2.3~\%$ of one
\omterm{def:b2:application-layer:resolver}{resolver}'s lookups reach the town's server; each \omterm{def:b2:application-layer:resolver}{resolver} has its own cache.

\textbf{10.} For up to the record's lifetime, an hour. The town lowers the lifetime (to a minute, say) at least an
hour before the move, then moves, then raises it again.

\textbf{11.} The records of an answer belong together (an alias and its target's address, a referral and its glue); a
part kept after another has expired could point to data the zone no longer holds.

\textbf{12.} One \omterm{def:b2:application-layer:resolver}{resolver} serves many machines, so its cache is shared; a machine waits one round trip to a nearby
\omterm{def:b2:application-layer:resolver}{resolver}, and needs no code to follow referrals.

\textbf{13.} \code{urllib}: 20; \code{http.client}: 1.

\textbf{14.} 19 handshakes of 80 ms: about 1.5 s.

\textbf{15.} A request's headers end with an empty line, and these requests have no body (a request with one says its
length). A response's body length is given by \code{Content-Length}.

\textbf{16.} Responses on an HTTP/1.1 connection come in the order of the requests, and nothing in a response says
which request it answers.

\textbf{17.} All four, which complete at 107 ms instead of 78 to 81: by 29, 28, 27 and 26 ms.

\textbf{18.} TCP delivers one byte stream in order, and holds back all bytes after a lost segment; it knows nothing of
HTTP's \omterm{def:b2:application-layer:streams}{streams}. QUIC orders the bytes of each stream separately.

\textbf{19.} 370 ms, of which 130 are the name.

\textbf{20.} 240 ms on a new connection, 80 ms on an open one.

\textbf{21.} One round trip per new connection: a cold first visit would take 450 ms.

\textbf{22.} Three TCP and three TLS handshakes, one per server; the 30 requests then cost round trips but no new
handshakes.

\textbf{23.} TLS hides the pages, the paths, the headers and the data sent. An observer still sees the addresses and
ports, the times and sizes of the records, the name lookups (in clear by default) and the site's name in the
ClientHello.

\textbf{24.} A first visit costs the name lookup (one round trip to the \omterm{def:b2:application-layer:resolver}{resolver} plus one per name server it must
ask) and three round trips: TCP, TLS 1.3 and the request. Here 370 ms; a revisit costs 240 ms on a new connection and
80 ms on an open one.
\end{solution}
